\documentclass[12pt,a4paper]{article}  % Déclare la classe du document.

\newcommand{\chaptermark}[2]{\og #2 \fg{} (#1)} % ligne uniquement pour éviter une erreur en cas d'article
\include{../1ere-preambule}

\usepackage{epic}   % extension et simplification de  l'environnement picture
\usepackage{graphicx}    % permet d'inclure des graphique externe dans un document
\usepackage{arydshln}   %permet d'avoir des lignes en pointillés

\newcommand{\lignes}[1]{
\setlength{\unitlength}{1mm}
\vskip 0,4cm\noindent
\multido{\i=0+1}{#1}{%
     \noindent
     \begin{picture}(150,5)
     \linethickness{0,005mm}
     \put(-5,0){\dashline{1}(175,5)(5,5)}
     \end{picture}\vskip 0,001cm
     }
\vskip 0cm
}

\rhead{DS n$^{\circ}$04 - 02/03/2023}

\begin{document}
\graphicspath{{../images/}}

\begin{center}
\fbox{
\begin{minipage}[]{15cm}
\begin{center}
\LARGE{\rule{0.5cm}{0pt}\rule[-0.25cm]{0pt}{0.9cm}
DS n$^{\circ}$04 - Produit scalaire
}
\end{center}
\end{minipage}
}

\end{center}


\begin{center}\textbf{2h - Calculatrice autorisée - Barème indicatif}\\
\textbf{Les élèves avec un tier-temps ne traitent pas les questions avec  le symbole \ding{46}}\\
Les résultats doivent être justifiés par des calculs (au moins 1 étape intermédiaire).
\end{center}

%\begin{center}\textbf{Les élèves avec un tier-temps ne traitent pas les questions avec  le symbole \ding{46}}\end{center}
%
%\begin{center}
%%On veillera à ENCADRER ses résultats et soigner sa copie.\\
%Les résultats doivent être justifiés par des calculs (au moins 1 étape intermédiaire).
%\end{center}

\vspace{-.3cm} 
\exerciceDS{ - ( 2 points )}
On considère 3 points $A, B$ et $C$. Dans chacun des cas suivants, calculer $\vv{A B} \cdot \vv{A C}$ ( on donnera la valeur exacte).
\begin{multicols}{2}
\begin{enumMet}
    \item $A B=4, A C=5$ et $\widehat{A B C}=\dfrac{\pi}{6}$
    \item $A B=6, A C=6$ et $\widehat{A B C}=\dfrac{\pi}{2}$
    \item $A B=3, A C=7$ et $\widehat{A B C}=\dfrac{\pi}{3}$
    \item $A B=4, A C=9$ et $\widehat{A B C}=\dfrac{\pi}{4}$
\end{enumMet}
\end{multicols}

\exerciceDS{ - ( 2 points )}
Dans les 2 cas ci dessous, calculer $\vv{A B} \cdot \vv{A C}$.
\begin{multicols}{2}
	\begin{enumMet}
		\item \ \\
		\newrgbcolor{ududff}{0.30196078431372547 0.30196078431372547 1.}
		\newrgbcolor{qqwuqq}{0. 0.39215686274509803 0.}
		\psset{xunit=0.9cm,yunit=0.9cm,algebraic=true,dimen=middle,dotstyle=o,dotsize=5pt 0,linewidth=1.6pt,arrowsize=3pt 2,arrowinset=0.25}
		\begin{pspicture*}(0.,0.5)(7.,6.5)
			\pspolygon[linewidth=2.pt,linecolor=qqwuqq,fillcolor=YellowGreen,fillstyle=solid,opacity=0.1](2.9376,5.1452)(3.3219,4.9655)(3.5016,5.3499)(3.1172,5.5295)
			\psline(1.,1.)(6.1218637165966525,4.125141927754758)
			\psline(6.1218637165966525,4.125141927754758)(2.2703662826677307,3.7177508178383416)
			\psline(1.,1.)(2.2703662826677307,3.7177508178383416)
			\psline[linestyle=dashed](2.2703662826677307,3.7177508178383416)(3.1172771377795505,5.529584696397235)
			\psline[linestyle=dashed](3.1172771377795505,5.529584696397235)(6.1218637165966525,4.125141927754758)
			\rput[bl](0.62,1.06){\ududff{$A$}}
			\rput[bl](2.,4.){\darkgray{$C$}}
			\rput[bl](2.84,5.84){\darkgray{$H$}}
			\rput[bl](6.2,4.28){\darkgray{$B$}}
			\rput[bl](3.72,2.3){$6$}
			\rput[bl](1.48,2.62){$3$}
			\rput[bl](2.5,4.82){$2$}
		\end{pspicture*}
		%\includegraphics[scale=1.2]{1.04.06-Img03}
		\item \ \\
		\newrgbcolor{ududff}{0.30196078431372547 0.30196078431372547 1.}
		\newrgbcolor{qqwuqq}{0. 0.39215686274509803 0.}
		\psset{xunit=1.0cm,yunit=1.0cm,algebraic=true,dimen=middle,dotstyle=o,dotsize=5pt 0,linewidth=1.3pt,arrowsize=3pt 2,arrowinset=0.25}
		\begin{pspicture*}(0.,0.2)(6.,4.)
			\pspolygon[linecolor=qqwuqq,fillcolor=YellowGreen,fillstyle=solid,opacity=0.1](2.8055522163904247,2.8369531245015627)(3.1081271803714507,2.5395503821270498)(3.4055299227459637,2.8421253461080758)(3.1029549587649377,3.1395280884825887)
			\psline(1.,1.)(4.989841373819047,1.284896142061883)
			\psline(4.989841373819047,1.284896142061883)(1.700984986254979,1.7131760294941962)
			\psline(1.,1.)(1.700984986254979,1.7131760294941962)
			\psline[linestyle=dashed](1.700984986254979,1.7131760294941962)(3.1029549587649377,3.1395280884825887)
			\psline[linestyle=dashed](3.1029549587649377,3.1395280884825887)(4.989841373819047,1.284896142061883)
				%\psdots[dotstyle=*,linecolor=ududff](1.,1.)
				\rput[bl](0.62,1.){\ududff{$B$}}
				%\psdots[dotsize=4pt 0,dotstyle=*,linecolor=darkgray](1.700984986254979,1.7131760294941962)
				\rput[bl](1.35,1.85){\darkgray{$A$}}
				%\psdots[dotsize=4pt 0,dotstyle=*,linecolor=darkgray](3.1029549587649377,3.1395280884825887)
				\rput[bl](2.82,3.46){\darkgray{$H$}}
				%\psdots[dotsize=4pt 0,dotstyle=*,linecolor=darkgray](4.989841373819047,1.284896142061883)
				\rput[bl](5.06,1.44){\darkgray{$C$}}
				\rput[bl](3.,0.8){$4$}
				\rput[bl](1.12,1.45){$1$}
				\rput[bl](2.1,2.5){$2$}
		\end{pspicture*}
		%\includegraphics[scale=1.2]{1.04.06-Img04}
	\end{enumMet}
\end{multicols}

%\vspace{-0.3cm}
%\exerciceDS{ - ( 1,5 points )}
%On considère un carré $A B C D$ de côté 5 cm. Calculer:
%\begin{multicols}{3}
%\begin{enumMet}
%    \item $\vv{A B} \cdot \vv{A C}$
%    \item $\vv{C A} \cdot \vv{C D}$
%    \item $\vv{D C} \cdot \vv{D A}$
%\end{enumMet}
%\end{multicols}

\vspace{-0.3cm}
\exerciceDS{ - ( 5 points )}

$\mathrm{ABC}$ est un triangle et $\mathrm{H}$ est le pied de la hauteur issue de $\mathrm{A}$. On a $\mathrm{AB}=5, \mathrm{BH}=3$ et $\mathrm{HC}=6$
\begin{multicols}{2}
\begin{enumMet}
    \item Calculer $\vv{B A} \cdot \vv{B C}$
    \item Calculer AH
    \item Calculer $\vv{A B} \cdot \vv{A H}$. On prendra si nécessaire $AH=4$.
    \item En déduire la mesure, en degré, de l'angle $\widehat{B A H}$. Arrondir au dixième.
\end{enumMet}
\columnbreak
\begin{center}
	\newrgbcolor{ududff}{0.30196078431372547 0.30196078431372547 1.}
	\newrgbcolor{qqwuqq}{0. 0.39215686274509803 0.}
	\psset{xunit=0.6cm,yunit=0.6cm,algebraic=true,dimen=middle,dotstyle=o,dotsize=5pt 0,linewidth=1.6pt,arrowsize=3pt 2,arrowinset=0.25}
	\begin{pspicture*}(1.,0.)(12.,6.)
		\pspolygon[linewidth=1.5pt,linecolor=qqwuqq,fillcolor=GreenYellow,fillstyle=solid,opacity=0.1](8.,1.4242640687119286)(7.575735931288071,1.4242640687119286)(7.575735931288071,1.)(8.,1.)
		\psline(2.,1.)(8.,1.)
		\psline(8.,1.)(11.,1.)
		\psline(2.,1.)(8.,5.)
		\psline(8.,5.)(11.,1.)
		\psline(8.,5.)(8.,1.)
		%\psdots[dotstyle=*,linecolor=ududff](2.,1.)
		\rput[bl](1.52,1.22){\ududff{$C$}}
		%\psdots[dotstyle=*,linecolor=ududff](8.,1.)
		\rput[bl](8.08,1.2){\ududff{$H$}}
		%\psdots[dotstyle=*,linecolor=ududff](11.,1.)
		\rput[bl](11.08,1.2){\ududff{$B$}}
		%\psdots[dotsize=4pt 0,dotstyle=*,linecolor=darkgray](8.,5.)
		\rput[bl](8.08,5.16){\darkgray{$A$}}
	\end{pspicture*}
    %\includegraphics[scale=0.6]{1.04.06-Img01}
\end{center}
\end{multicols}

%\vspace{-0.3cm}
%\exerciceDS{ - ( 2 points )}
%
%Soit un losange $A B C D$ de centre $O$ tel que $A C=4$ et $B D=7$.
%%\begin{multicols}{2}
%\begin{enumMet}
%    \item Calculer $\vv{B C} \cdot \vv{B D}$.
%    \item Calculer $\vv{D C} \cdot \vv{D O}$.
%\end{enumMet}
%%\end{multicols}
%
%\vspace{-0.3cm}
%\exerciceDS{ - ( 3 points )}
%On considère les points $A(-1 ; 2), B(4 ; 3)$ et $C(1 ;-5)$
%\begin{enumMet}
%    \item Calculer les coordonnées des vecteurs $\vv{A B}, \vv{A C}$ et $\vv{B C}$
%    \item Calculer $\vv{A B} \cdot \vv{A C}, \vv{A B} \cdot \vv{B C}$ et $\vv{A C} \cdot \vv{B C}$
%\end{enumMet}
%\newpage

\newpage
\vspace{-0.3cm}
\exerciceDS{ - ( 3 points )}
Soit les points $A(2 ; 7), B(6 ; 5)$ et $C(4 ; 1)$.
%\begin{multicols}{2}
\begin{enumMet}
    \item Calculer le produit scalaire $\vv{B A} \cdot \vv{B C}$
    \item Que peut-on en conclure pour le triangle $A B C$ ?
\end{enumMet}
%\end{multicols}


\exerciceDS{ - ( 3 points )}

Dans le repère orthonormé ci-dessous, on donne le triangle $A B C$.\\
\begin{center}
	\newrgbcolor{ududff}{0.30196078431372547 0.30196078431372547 1.}
	\psset{xunit=0.7cm,yunit=0.7cm,algebraic=true,dimen=middle,dotstyle=o,dotsize=10pt 0,linewidth=1.6pt,arrowsize=3pt 2,arrowinset=0.25}
	\begin{pspicture*}(-1.,-0.8)(7.5,5.8)
		\psgrid[gridwidth=0.6pt,gridcolor=gray,subgriddiv=0,gridlabels=0](0,0)(-1.,-1)(8,6)
		\psaxes[labelFontSize=\scriptstyle,xAxis=true,yAxis=true,Dx=1.,Dy=1.,ticksize=-2pt 0,subticks=2]{->}(0,0)(-1.,-0.8)(7.5,5.8)
		\psline(1.,2.)(3.,5.)
		\psline(3.,5.)(7.,1.)
		\psline(7.,1.)(1.,2.)
		\psdots[dotstyle=x,linecolor=ududff](1.,2.)
		\rput[bl](0.45,2.12){\ududff{$A$}}
		\psdots[dotstyle=x,linecolor=ududff](7.,1.)
		\rput[bl](7.08,1.2){\ududff{$B$}}
		\psdots[dotstyle=x,linecolor=ududff](3.,5.)
		\rput[bl](3.08,5.2){\ududff{$C$}}
	\end{pspicture*}
    %\includegraphics[scale=0.95]{1.04.06-Img02}
\end{center}
%\begin{enumMet}
%    \item 
    \begin{enumMet}
        \item calculer $\vv{A B} \cdot \vv{A C}$.
        \item En déduire la mesure, en degré, de l'angle $\widehat{B A C}$. Arrondir au dixième.
    \end{enumMet}
%    \item \begin{enumMet}
%        \item Dans ce triangle, H est le pied de la hauteur issue du sommet $C$.
%Déterminer la longueur $\mathrm{CH}$. Arrondir au dixième.
%        \item En déduire une valeur approchée de l'aire du triangle $A B C$.
%    \end{enumMet}
%\end{enumMet}

\exerciceDS{ - ( 2 points )}
Soient deux vecteurs $\vv{u}$ et $\vv{v}$ tels que $\|\vv{u}\|=2,\|\vv{v}\|=4$ et $(\vv{u} ; \vv{v})=\dfrac{\pi}{3}$.
\begin{enumMet}
	\item Calculer $\vv{u}\cdot\vv{v}$
	\item Calculer $(\vv{u}-3 \vv{v}) \cdot(5 \vv{u}+\vv{v})$
\end{enumMet}


\exerciceDS{ - ( 4 points )}

$A B C D$ est un rectangle tel que $A B=4$ et $B C=3$. Le triangle $A B F$ est équilatéral et le triangle $B C E$ est rectangle isocèle en $C$. H est le milieu de $[A B]$. Calculer les produits scalaires suivants : 
\begin{multicols}{2}
	\begin{multicols}{2}
	\begin{enumMet}
		\item $ \vv{A B} \cdot  \vv{A C} \quad$ 
		\item $ \vv{A B} \cdot  \vv{A H}$
		%\item $ \vv{B A} \cdot  \vv{B F}$
		\item $ \vv{A D} \cdot  \vv{C E} \quad$ 
		\item $ \vv{B A} \cdot  \vv{A F}$
		\item $ \vv{C A} \cdot  \vv{C E}$
		\item $ \vv{C E} \cdot  \vv{B A} \quad$
		\item $ \vv{A F} \cdot  \vv{D E}$
		\item $ \vv{B E} \cdot  \vv{D A}$
	\end{enumMet}
	\end{multicols}

	\columnbreak
	
	\newrgbcolor{ududff}{0.30196078431372547 0.30196078431372547 1.}
	\newrgbcolor{xdxdff}{0.49019607843137253 0.49019607843137253 1.}
	\newrgbcolor{qqwuqq}{0. 0.39215686274509803 0.}
	\psset{xunit=0.8cm,yunit=0.8cm,algebraic=true,dimen=middle,dotstyle=o,dotsize=10pt 0,linewidth=1.5pt,arrowsize=3pt 2,arrowinset=0.25}
	\begin{pspicture*}(-1.,-0.6)(8.,7.)
		\pspolygon[linecolor=qqwuqq,fillcolor=YellowGreen,fillstyle=solid,opacity=0.1](4.424264068711929,0.)(4.424264068711929,0.4242640687119286)(4.,0.4242640687119286)(4.,0.)
		\psline(4.,3.)(4.,0.)
		\psline(4.12,1.6)(3.88,1.6)
		\psline(4.12,1.5)(3.88,1.5)
		\psline(4.12,1.4)(3.88,1.4)
		\psline(4.,0.)(0.,0.)
		\psline(0.,0.)(0.,3.)
		\psline(4.,3.)(7.,0.)
		\psline(7.,0.)(4.,0.)
		\psline(5.6,-0.12)(5.6,0.12)
		\psline(5.5,-0.12)(5.5,0.12)
		\psline(5.4,-0.12)(5.4,0.12)
		\psline(0.,3.)(2.,6.464101615137754)
		\psline(2.,6.464101615137754)(4.,3.)
		\psline(0.,3.)(2.,3.)
		\psline(0.95,3.12)(0.95,2.88)
		\psline(1.05,3.12)(1.05,2.88)
		\psline(2.,3.)(4.,3.)
		\psline(2.95,3.12)(2.95,2.88)
		\psline(3.05,3.12)(3.05,2.88)
			\rput[bl](-0.54,3.2){\ududff{$A$}}
			\rput[bl](4.08,3.2){\ududff{$B$}}
			\rput[bl](3.54,0.12){\xdxdff{$C$}}
			\rput[bl](-0.56,0.16){\darkgray{$D$}}
			\rput[bl](7.08,0.2){\xdxdff{$E$}}
			\rput[bl](2.08,6.62){\darkgray{$F$}}
			\psdots[dotstyle=x,linecolor=xdxdff](2.,3.)
			\rput[bl](2.08,3.2){\xdxdff{$H$}}
	\end{pspicture*}
\end{multicols}



\end{document}

